\subsection{The graded module associated with a filtered module}

Let G be a commutative group (written additively) and \((G_{n})\) a filtration on G. Let us write

\[
\begin{array}{l} \operatorname{gr} _ {n} (\mathbf {G}) = \mathbf {G} _ {n} / \mathbf {G} _ {n + 1} \quad \text { for } \quad n \in \mathbf {Z} \\ \operatorname{gr} (\mathbf {G}) = \bigoplus_ {n \in \mathbf {Z}} \operatorname{gr} _ {n} (\mathbf {G}). \end{array}\tag{8}
\]

The commutative group \(\operatorname{gr}(\mathbf{G})\) is then a graded group of type \(\mathbf{Z}\) , called the graded group associated with the filtered group \(\mathbf{G}\) , the homogeneous elements of degree \(n\) of \(\operatorname{gr}(\mathbf{G})\) being those of \(\operatorname{gr}_n(\mathbf{G})\) .

Now let A be a filtered ring, \((\mathbf{A}_{n})\) its filtration, E a filtered A-module and \((\mathbf{E}_{n})\) its filtration. For all \(p \in Z\) , \(q \in Z\) , a mapping

\[
\operatorname{gr} _ {p} (\mathrm{A}) \times \operatorname{gr} _ {q} (\mathrm{E}) \rightarrow \operatorname{gr} _ {p + q} (\mathrm{E})\tag{9}
\]

is defined as follows: given \(\alpha \in \mathrm{gr}_p(\mathbf{A})\) , \(\xi \in \mathrm{gr}_q(\mathbf{E})\) , two representatives \(a, a'\) of \(\alpha\) and two representatives \(x, x'\) of \(\xi\) , \(ax \in \mathrm{E}_{p + q}\) , \(a'x' \in \mathrm{E}_{p + q}\) and \(ax \equiv a'x'\) (mod. \(\mathrm{E}_{p + q + 1}\) ), for

\[
a x - a ^ {\prime} x ^ {\prime} = (a - a ^ {\prime}) x + a ^ {\prime} (x - x ^ {\prime})
\]

and \(a - a' \in \mathbf{A}_{p+1}\) and \(x - x' \in \mathbf{E}_{q+1}\) and hence our assertion follows from formula (3) of no. 1. We may therefore denote by \(\alpha\xi\) the class in

\[
\mathbf {E} _ {p + 1} / \mathbf {E} _ {p + q + 1} = \operatorname{gr} _ {p + q} (\mathbf {E})
\]

of the product ax of any representative \(a \in \alpha\) and any representative \(x \in \xi\) . It is immediate that the mapping (9) is Z-bilinear; by linearity, we derive a Z-bilinear mapping

\[
\operatorname{gr} (\mathrm{A}) \times \operatorname{gr} (\mathrm{E}) \rightarrow \operatorname{gr} (\mathrm{E}).\tag{10}
\]

If this definition is first applied to the case \(E = A_{s}\) , the mapping (10) is an internal law of composition on \(\mathrm{gr}(A)\) , which it is immediately verified is associative and has an identity element which is the canonical image in \(\mathrm{gr}_{0}(A)\) of the unit element of A; it therefore defines on \(\mathrm{gr}(A)\) a ring structure and the graduation \((\mathrm{gr}_{n}(A))_{n\in\mathbf{Z}}\) is by definition compatible with this structure. The graded ring \(\mathrm{gr}(A)\) (of type Z) thus defined is called the graded ring associated with the filtered ring A; it is obviously commutative if A is commutative; \(\mathrm{gr}_{0}(A)\) is a subring of \(\mathrm{gr}(A)\) . The mapping (10) is on the other hand a \(\mathrm{gr}(A)\) -module external law on \(\mathrm{gr}(E)\) , the module axioms being trivially satisfied, and the graduation \((\mathrm{gr}_{n}(E))_{n\in\mathbf{Z}}\) on \(\mathrm{gr}(E)\) is obviously compatible with this module structure. The graded \(\mathrm{gr}(A)\) -module \(\mathrm{gr}(E)\) (of type Z) thus defined is called the graded module associated with the filtered A-module E.

\textbf{Examples}\par

\begin{enumerate}
\def\labelenumi{(\arabic{enumi})}
\item
  Let A be a commutative ring and t an element of A which is not a divisor of 0. Let us give A the \((t)\) -adic filtration (no. 1, Example 3). Then the associated graded ring \(\mathrm{gr}(\mathbf{A})\) is canonically isomorphic to the polynomial ring \((\mathbf{A}/(t))[\mathbf{X}]\) . For \(\mathrm{gr}_{n}(\mathbf{A}) = 0\) for n \textless{} 0 and by definition the ring \(\mathrm{gr}_{0}(\mathbf{A})\) is the ring \(\mathbf{A}/(t)\) . We now note that by virtue of the hypothesis on t the relation \(at^{n} \equiv 0 \pmod{t^{n+1}}\) is equivalent to \(a \equiv 0 \pmod{t}\) ; if \(\tau\) is the canonical image of t in \(\mathrm{gr}_{1}(\mathbf{A})\) , every element of \(\mathrm{gr}_{n}(\mathbf{A})\) may then be written uniquely in the form \(\alpha\tau^{n}\) where \(\alpha \in \mathrm{gr}_{0}(\mathbf{A})\) ; whence our assertion.
\item
  Let \(\mathbf{K}\) be a commutative ring, A the ring of formal power series

\[
\mathrm{K} \left[ \left[ \mathrm{X} _ {1}, \dots , \mathrm{X} _ {r} \right] \right]
\]

(Algebra, Chapter IV, § 5) and m the ideal of A whose elements are the formal power series with no constant term. Let us give A the m-adic filtration (no. 1, Example 3); if \(M_{1}, \ldots, M_{s}\) are the distinct monomials in \(X_{1}, \ldots, X_{r}\) of total degree n - 1, clearly every formal power series u of total order \(\omega(u) \geqslant n\)

(loc. cit., no. 2) may be written as \(\sum_{k=1} u_k M_k\) , where the \(u_k\) belong to \(m\) ; it is seen that \(m^n\) is the set of formal power series \(u\) such that \(\omega(u) \geqslant n\) , which shows that \(\omega\) is the order function for the m-adic filtration. Then clearly, for every formal power series \(u \in m^n\) , there exists a unique homogeneous polynomial of degree \(n\) in the \(X_i\) which is congruent to \(u\) mod. \(m^{n+1}\) , namely the sum of terms of degree \(n\) of \(u\) ; we conclude that \(\text{gr}(A)\) is canonically isomorphic to the polynomial ring \(K[X_1, \ldots, X_f]\) .

\item
  More generally, let A be a commutative ring, b an ideal of A and A be given the b-adic filtration. If we write \(\mathbf{B} = \mathbf{gr}_{0}(\mathbf{A})\) , \(\mathbf{F} = \mathbf{gr}_{1}(\mathbf{A}) = \mathbf{b}/\mathbf{b}^{2}\) , we know (Algebra, Chapter III,) that the identity mapping of the B-module F onto itself can be extended uniquely to a homomorphism u from the symmetric algebra \(\mathbf{S}(\mathbf{F})\) of F to the B-algebra \(\mathbf{gr}(\mathbf{A})\) ; it follows from the definition of \(\mathbf{gr}(\mathbf{A})\) that u is a surjective homomorphism of graded algebras; for \(n \geqslant 1\) , every element of \(\mathbf{gr}_{n}(\mathbf{A})\) is a sum of classes mod \(b^{n+1}\) of elements of the form \(y = x_{1}x_{2}\ldots x_{n}\) , where \(x_{i} \in b (1 \leqslant i \leqslant n)\) ; if \(\xi_{i}\) is the class of \(x_{i} \mod b^{2}\) , clearly the class of y mod. \(b^{n+1}\) is the element \(u(\xi_{1})\ldots u(\xi_{n})\) , whence our assertion. In particular, every system of generators of the B-module F is a system of generators of the B-algebra \(\mathbf{gr}(\mathbf{A})\) .
\end{enumerate}

If now E is an A-module and E is given the b-adic filtration, it is seen similarly that the graded gr(A)-module gr(E) is generated by gr₀(E) = E/bE. To be precise, the restriction φ to gr(A) × gr₀(E) of the external law on the gr(A)-module gr(E) is a Z-bilinear mapping of gr(A) × gr₀(E) to gr(E); moreover gr(A) is a (gr₀(A), gr₀(A))-bimodule and gr₀(E) a gr₀(A)-module; it is immediately verified that, for α ∈ gr(A), α₀ ∈ gr₀(A), ξ ∈ gr₀(E),

\[
\phi (\alpha \alpha_ {0}, \xi) = \phi (\alpha , \alpha_ {0} \xi)
\]

and hence \(\phi\) defines a surjective \(\mathbf{gr}_0(\mathbf{A})\) -linear mapping

\[
\gamma_ {\mathrm{E}} \colon \operatorname{gr} (\mathrm{A}) \otimes_ {\mathbf {g r} _ {0} (\mathrm{A})} \operatorname{gr} _ {0} (\mathrm{E}) \rightarrow \operatorname{gr} (\mathrm{E})\tag{11}
\]

which is called canonical.

* (4) Let K be a commutative ring, g a Lie algebra over K and U the enveloping algebra of g. An increasing filtration \((\mathrm{U}_{n})_{n\in\mathbf{Z}}\) is defined on U by taking \(U_{n}=\{0\}\) for n\textless0 and denoting by \(U_{n}\) for \(n\geqslant0\) the set of elements of U which can be expressed as a sum of products of at most n elements of g; then

\(U_{0} = K\) and \(\text{gr}(U)\) is a commutative K-algebra (Lie Groups and Lie Algebras, Chapter I, § 2, no. 6). The canonical mapping of \(g\) to \(\text{gr}_{1}(U) = U_{1}/U_{0}\) can be extended uniquely to a homomorphism \(h\) of the symmetric algebra \(S(g)\) of the K-module \(g\) to the K-algebra \(\text{gr}(U)\) ; the homomorphism \(h\) is surjective and, if the K-module \(g\) is free, \(h\) is bijective (loc. cit., no. 7, Theorem 1).

\begin{enumerate}
\def\labelenumi{(\arabic{enumi})}
\setcounter{enumi}{4}
\tightlist
\item
  Let A be a graded ring of type Z and E a graded A-module of type Z; let \(A_{(i)}\) (resp. \(E_{(i)}\) ) be the subgroup of homogeneous elements of degree i of A (resp. E). Let A and E be given the filtrations associated with their graduations (no. 1, Example 1) and let \(A'\) and \(E'\) denote the filtered ring and filtered \(A'\) -module thus obtained. Then it is immediate that the Z-linear mapping \(A \to \text{gr}(A')\) which maps an element of \(A_{(n)}\) to its canonical image in

\[
\operatorname{gr} _ {n} (\mathbf {A}) = \left(\bigoplus_ {i \geqslant n} \mathbf {A} _ {(i)}\right) / \left(\bigoplus_ {i \geqslant n + 1} \mathbf {A} _ {(i)}\right)
\]

is a graded ring isomorphism. A canonical graded A-module isomorphism \(\mathbf{E} \to \mathrm{gr}(\mathbf{E}')\) is defined similarly.
\end{enumerate}

PROPOSITION 1. Let A be a filtered ring, \((\mathbf{A}_n)_{n\in \mathbf{Z}}\) its filtration and \(v\) its order function. Suppose that \(\operatorname {gr}(\mathbf{A})\) is a ring with no divisor of zero. Then, for every ordered pair of elements \(a,b\) of the ring \(\mathbf{B} = \bigcup_{n\in \mathbf{Z}}\mathbf{A}_n,v(ab) = v(a) + v(b)\) .

As \(n = \bigcap_{n \in \mathbf{Z}} A_n\) is a two-sided ideal of the ring \(B\) , the formula holds if \(v(a)\) or \(v(b)\) is equal to \(+\infty\) . If not, \(v(a) = r\) and \(v(b) = s\) are integers; the classes \(\alpha\) of \(a \bmod A_{r+1}\) and \(\beta\) of \(b \bmod A_{s+1}\) are \(\neq 0\) by definition, whence by hypothesis \(\alpha \beta \neq 0\) in \(\operatorname{gr}(A)\) and therefore \(ab \notin A_{r+s+1}\) ; as \(ab \in A_{r+s}\) ,

\[
v (a b) = v (a) + v (b).
\]

COROLLARY. Let A be a filtered ring and \((\mathbf{A}_{n})_{n\in\mathbf{Z}}\) its filtration; let us set \(B=\bigcup_{n\in Z}A_{n}\) , \(n=\bigcap_{n\in Z}A_{n}\) . If the ring gr(A) has no divisors of zero, neither has the ring B/n.

If \(a\) and \(b\) are elements of \(B\) not belonging to \(\mathfrak{n}\) , then \(v(a) \neq +\infty\) and \(v(b) \neq +\infty\) , whence \(v(ab) \neq +\infty\) and therefore \(ab \notin \mathfrak{n}\) .

Note that the ring A can be an integral domain and the filtration \((A_{n})\) exhaustive and separated without gr(A) being an integral domain (Exercise 2).

Remark. Let G be a group which is not necessarily commutative with a filtration \((\mathrm{G}_{n})_{n\in\mathbf{Z}}\) such that \(G_{n+1}\) is normal in \(G_{n}\) for all \(n\in Z\) ; again let \(\mathrm{gr}_{n}(\mathrm{G})=\mathrm{G}_{n}/\mathrm{G}_{n+1}\) . The restricted product of the family \((\mathrm{gr}_{n}(\mathrm{G}))_{n\in\mathbf{Z}}\) , that is the subgroup of the product \(\prod_{n\in\mathbf{Z}}\mathrm{gr}_{n}(\mathrm{G})\) consisting of the elements \((\xi_{n})\) all of whose components, except at most a finite number, are equal to the identity element, is also called the graded group associated with G and denoted by \(\mathrm{gr}(\mathrm{G})\) .

